{\large\textbf{The Expanding Universe}\hfill\textbf{(Total Marks : 20)}}

The most outstanding fact in cosmology is that our universe is expanding. Space is continuously created as time lapses. The expansion of space indicates that, when the universe expands, the distance between objects in our universe also expands. It is convenient to use “comoving” coordinate system $\vec{r} = (x, y, z)$ to label points in our expanding universe, in which the coordinate distance $\Delta r = |\vec{r}_2 - \vec{r}_1| = \sqrt{(x_2 - x_1)^2 + (y_2 - y_1)^2 + (z_2 - z_1)^2}$ between objects 1 and 2 does not change. (Here we assume no peculiar motion, i.e.\ no additional motion of those objects other than the motion following the expansion of the universe.) The situation is illustrated in the figure below (the figure has two space dimensions, but our universe actually has three space dimensions).

\begin{center}
\includegraphics[width=0.75\linewidth]{figures/APhO_2016_Q2_fig1.png}
\end{center}

The modern theory of cosmology is built upon Einstein’s general relativity. However, under proper assumptions, a simplified understanding under the framework of Newton’s theory of gravity is also possible. In the following questions, we shall work in the framework of Newton’s gravity.

To measure the physical distance, a “scale factor” $a(t)$ is introduced such that the physical distance $\Delta r_\mathrm{p}$ between the comoving points $\vec{r}_1$ and $\vec{r}_2$ is
\[
\Delta r_\mathrm{p} = a(t) \Delta r,
\]
The expansion of the universe implies that $a(t)$ is an increasing function of time.

On large scales -- scales much larger than galaxies and their clusters -- our universe is approximately homogeneous and isotropic. So let us consider a toy model of our universe, which is filled with uniformly distributed particles. There are so many particles, such that we model them as a continuous fluid. Furthermore, we assume the number of particles is conserved.

Currently, our universe is dominated by non-relativistic matter, whose kinetic energy is negligible compared to its mass energy. Let $\rho_\mathrm{m}(t)$ be the physical energy density (i.e.\ energyper unit physical volume, which is dominated by mass energy for non-relativistic matter and the gravitational potential energy is not counted as part of the “physical energy density”) of non-relativistic matter at time $t$. We use $t_0$ to denote the present time.

\noindent\begin{tabular}{|c|p{0.7\linewidth}|c|}
\hline
\textbf{A} & Derive the expression of $\rho_\mathrm{m}(t)$ at time $t$ in terms of $a(t)$, $a(t_0)$ and $\rho_\mathrm{m}(t_0)$. & \textbf{2 points} \\
\hline
\end{tabular}

Besides non-relativistic matter, there is also a small amount of radiation in our current universe, which is made of massless particles, for example, photons. The physical wavelength of massless particles increases with the universe expansion as $\lambda_\mathrm{p} \propto a(t)$. Let the physical energy density of radiation be $\rho_\mathrm{r}(t)$.

\noindent\begin{tabular}{|c|p{0.7\linewidth}|c|}
\hline
\textbf{B} & Derive the physical energy density for radiation $\rho_\mathrm{r}(t)$ at time $t$ in terms of $a(t)$, $a(t_0)$ and $\rho_\mathrm{r}(t_0)$. & \textbf{2 points} \\
\hline
\end{tabular}

Consider a gas of non-interacting photons which has thermal equilibrium distribution. In this situation, the temperature of the photon depends on time as $T(t) \propto [a(t)]^\gamma$.

\noindent\begin{tabular}{|c|p{0.7\linewidth}|c|}
\hline
\textbf{C} & Calculate the numerical value of $\gamma$. & \textbf{2 points} \\
\hline
\end{tabular}

Consider the thermodynamics of one type of non-interacting particle X. Note that the space expansion is slow enough and thermally isolated such that the entropy of X is a constant in time. Let the physical energy density of X be $\rho_\mathrm{X}(t)$, which includes mass energy and internal energy. Let the physical pressure be $p_\mathrm{X}(t)$.

\noindent\begin{tabular}{|c|p{0.7\linewidth}|c|}
\hline
\textbf{D} & Derive $\mathrm{d}\rho_\mathrm{X}(t)/\mathrm{d}t$ in terms of $a(t)$, $\mathrm{d}a(t)/\mathrm{d}t$,$\rho_\mathrm{X}(t)$,and $p_\mathrm{X}(t)$. & \textbf{4 points} \\
\hline
\end{tabular}

Consider a star S. At the present time $t_0$, the star is at a physical distance $r_\mathrm{p} = a(t_0)r$ away from us, where $r$ is the comoving distance. Here we ignore the peculiar motion, i.e.\ assume that both the star and us just follow the expansion of the universe without additional motion.

The star is emitting energy in the form of light at power $P_\mathrm{e}$, which is isotropic in every direction. We use a telescope to observe its starlight. For simplicity, assume the telescope can observe all frequencies of light with 100\% efficiency. Let the area of the telescope lens be $A$.

\noindent\begin{tabular}{|c|p{0.7\linewidth}|c|}
\hline
\textbf{E} & Derive the power received by the telescope $P_\mathrm{r}$ from the star S, as a function of $r, A, P_\mathrm{e}$, the scale factor $a(t_\mathrm{e})$ at the starlight emission time $t_\mathrm{e}$, and the present (i.e.\ at the observation time) scale factor $a(t_0)$. & \textbf{4 points} \\
\hline
\end{tabular}

If there were no gravity, the expansion speed of the universe should be a constant. In Newton’s framework, this can be understood as that, without force, matter just moves away from each other with constant speed and thus $\mathrm{d}a(t)/\mathrm{d}t$ is a constant depending on the initial condition.

Let us now consider how Newton’s gravity affects the scale factor $a(t)$, in a universe filled with non-relativistic matter in a homogeneous and isotropic way.

\begin{center}
\includegraphics[width=0.38\linewidth]{figures/APhO_2016_Q2_fig2.png}
\end{center}

As illustrated in the above figure, let us assume C is the center of our universe (this assumption can be removed in Einstein’s general relativity, which is beyond the scope of this question). We slice matter into thin shells around C. Let us focus on one thin shell (the sphere in the above figure) whose comoving distance from the center is $r$ (recall that this comoving distance is a constant in time).

\noindent\begin{tabular}{|c|p{0.7\linewidth}|c|}
\hline
\textbf{F} & Use the motion of the shell to find a relation between $\mathrm{d}a(t)/\mathrm{d}t$, $a(t)$ and the density of mass energy $\rho(t)$. (In the final relation, if you encounter a constant depending on the initial condition, it can be kept as it is.) & \textbf{5 points} \\
\hline
\end{tabular}

\noindent\begin{tabular}{|c|p{0.7\linewidth}|c|}
\hline
\textbf{G} & Based on the model described in Part (F), is the expansion of the universe (a) accelerating or (b) decelerating? Choose from (a) or (b). & \textbf{1 points} \\
\hline
\end{tabular}

For your information, in 1998, a new type of energy component of our universe is discovered. It actually changes the conclusion in Part (G).
